\chapter{Теоретическая часть}

\section{Обзор предметной области}

Текст раздела. Ссылка на рисунок~\ref{fig:example} и таблицу~\ref{tab:example}.
Пример формулы:
\begin{equation}
  E = mc^2,
  \label{eq:einstein}
\end{equation}
где $E$ --- энергия, $m$ --- масса, $c$ --- скорость света.

\begin{figure}[H]
  \centering
  % \includegraphics[width=0.7\textwidth]{имя_файла}  % картинки клади в images/
  \fbox{\parbox[c][4cm][c]{0.6\textwidth}{\centering Здесь будет рисунок}}
  \caption{Пример рисунка}
  \label{fig:example}
\end{figure}

\begin{table}[H]
  \caption{Пример таблицы}
  \label{tab:example}
  \centering
  \begin{tabular}{|l|c|c|}
    \hline
    Параметр & Значение 1 & Значение 2 \\ \hline
    A & 1 & 2 \\ \hline
    B & 3 & 4 \\ \hline
  \end{tabular}
\end{table}

\section{Постановка задачи}

Текст раздела~\cite{gost7322017}.
